\documentclass[11pt]{article}

\usepackage[T1]{fontenc}
\usepackage[utf8]{inputenc}
\usepackage{lmodern}
\usepackage{amsmath,amssymb,amsthm,mathtools}
\usepackage{booktabs}
\usepackage{enumitem}
\usepackage[margin=1in]{geometry}
\usepackage[colorlinks=true,linkcolor=blue,citecolor=blue,urlcolor=blue]{hyperref}
\usepackage[nameinlink,noabbrev]{cleveref}
\usepackage{microtype}

\newtheorem{theorem}{Theorem}[section]
\newtheorem{lemma}[theorem]{Lemma}
\newtheorem{proposition}[theorem]{Proposition}
\newtheorem{corollary}[theorem]{Corollary}
\theoremstyle{definition}
\newtheorem{definition}[theorem]{Definition}
\theoremstyle{remark}
\newtheorem{remark}[theorem]{Remark}

\newcommand{\fDelta}{f_{\Delta}}
\newcommand{\deficiency}{\operatorname{def}}
\newcommand{\N}{N}
\newcommand{\Z}{\mathbb Z}
\newcommand{\floor}[1]{\left\lfloor #1\right\rfloor}

\title{An Exact Extremal Formula for Triangle-Free Graphs\\
with Bounded Degree and Matching Number}
\author{Erdem Banak\\[2mm]
\small Research draft, not peer reviewed}
\date{October 2026}

\begin{document}

\maketitle

\begin{abstract}
We determine the maximum number of edges in a triangle-free graph with
prescribed upper bounds on maximum degree and matching number. The result
is an explicit formula for all degrees $d\ge2$ and matching bounds $m\ge0$,
together with constructions attaining it. We combine the component reduction
of Ahanjideh, Ekim, and Y\i ld\i z and the construction family of Banak, Ekim,
and Ta\c{s}k\i n with a new fixed-order bound. Its proof extracts a saturated
five-cycle, forces a homomorphism to $C_5$ under a total-deficiency hypothesis,
and uses an odd-excess imbalance to bound the missing edges. The same argument
determines the remaining odd-degree thresholds. An elementary packing
argument then yields the general formula, resolving the two conjectures.
The new fixed-order argument, including the shortest-odd-cycle step, and the
packing normalization are checked in Lean. The final deductions are
conditional on explicit published inputs.
\end{abstract}

\section{Introduction}

For integers $d\ge2$ and $m\ge0$, let $\fDelta(d,m)$ be the maximum number of
edges in a finite simple triangle-free graph $G$ satisfying
\[
  \Delta(G)\le d,\qquad \nu(G)\le m,
\]
where $\nu(G)$ denotes matching number. The main theorem gives this maximum
without an implicitly defined threshold parameter.

\begin{theorem}[General extremal formula]\label{thm:main}
For $d\ge2$, define
\begin{equation}\label{eq:explicit-threshold}
  z_d=
  \begin{cases}
    \floor{5d/4},&d\text{ even},\\
    3(d-1)/2,&d\in\{3,5,7\},\\
    \floor{5(d+1)/4},&d\ge9\text{ odd}.
  \end{cases}
\end{equation}
For every $m\ge0$, write $m=kz_d+r$, where $0\le r<z_d$. Then
\begin{equation}\label{eq:main-formula}
  \fDelta(d,m)=dm+k\floor{d/2}+\max\{0,r-d+1\}.
\end{equation}
The value is attained by a disjoint union of threshold components, at most
one additional non-star component, and $d$-stars.
\end{theorem}

Ahanjideh, Ekim, and Y\i ld\i z~\cite{AEY2022} proved the formula in several
parameter ranges and established the component reduction on which we build.
They introduced $Z(d)$, the least matching number of a triangle-free
factor-critical graph that is $d$-regular for even $d$, or has one vertex of
degree $d-1$ and all others of degree $d$ for odd $d$. Banak, Ekim, and
Ta\c{s}k\i n~\cite{BET2023} extended the computational results and supplied
the intermediate construction family. Their Conjectures~2.1 and~2.2 assert
the intermediate component value and the unrestricted extremal formula.

We integrate the short component and endpoint proofs from~\cite{AEY2022},
with attribution, into the argument below. The new step is the fixed-order
upper bound: a graph exceeding the proposed component value has too little
total deficiency to avoid a saturated five-cycle. The cycle excess
$q=2|V(G)|-5d$ then governs both the extension to a $C_5$-coloring and the
final missing-edge count. This proves the intermediate bound and identifies
$Z(d)=z_d$; a single packing argument completes \cref{thm:main}.

Thus the previously conjectured formulas are consequences of the general
theorem and its component bound. The theorem gives the extremal value and
examples attaining it; it does not classify every extremal graph.

\section{Reduction to factor-critical components}\label{sec:reduction}

We omit isolated vertices. The maximum defining $\fDelta(d,m)$ exists:
the endpoints of a maximal matching form a vertex cover of size at most
$2m$, so every admissible graph has at most $2md$ edges. Among graphs attaining
the maximum, choose one with as many $d$-star components as possible.
Call such a graph \emph{normalized}. It has matching number exactly $m$,
since otherwise adjoining a $d$-star would increase its edge count.
Below, a non-star component means one not isomorphic to $K_{1,d}$.

We use Gallai's lemma: a connected graph whose matching number is unchanged
by deletion of any vertex is factor-critical; see~\cite[Lemma~2.2]{AEY2022}.
The following reduction is the argument of~\cite[Lemma~2.4]{AEY2022}.

\begin{lemma}\label{lem:star-reduction}
Every non-star component of a normalized extremal graph is factor-critical.
It has $2j+1$ vertices, where $j$ is its matching number, and is itself
extremal for the bounds $(d,j)$.
\end{lemma}

\begin{proof}
If a non-star component $H$ were not factor-critical, Gallai's lemma would
give a vertex $v$ with $\nu(H-v)\le\nu(H)-1$. Replace $H$ by $H-v$ and a
$d$-star. This preserves admissibility and loses at most $d$ edges while
adding $d$, so the resulting graph is still extremal. It has more $d$-stars,
a contradiction. Factor-criticality gives the asserted order. Finally,
replacing $H$ by any admissible graph with more edges would contradict
extremality of the whole graph.
\end{proof}

The next bound is~\cite[Lemma~3.1]{AEY2022}; its proof uses Mantel's theorem.

\begin{lemma}\label{lem:fc-order-bound}
A connected triangle-free factor-critical graph $H$ with matching number
$j$ has at most $j^2+1$ edges.
\end{lemma}

\begin{proof}
The one-vertex case is immediate. Otherwise $H$ is nonbipartite and has
$2j+1$ vertices. Choose a shortest odd cycle of length $2s+1$, with $s\ge2$.
It is induced, and every outside vertex has at most $s$ neighbors on it.
The remaining $2(j-s)$ vertices induce at most $(j-s)^2$ edges by Mantel's
theorem. Consequently
\[
  e(H)\le(2s+1)+2s(j-s)+(j-s)^2
        =j^2+2-(s-1)^2\le j^2+1.
\]
\end{proof}

\begin{corollary}\label{lem:component-lower}
Every non-star component of a normalized extremal graph has matching number
at least $d$.
\end{corollary}

\begin{proof}
Here $j\ge2$: the case $j=0$ is an omitted isolated vertex, while the only
factor-critical graph on three vertices is $K_3$, which is not triangle-free.
If $j<d$, then
$j^2+1<dj$, and replacing it by $j$ copies of a $d$-star increases the edge
count, contradicting \cref{lem:fc-order-bound} and extremality.
This is the replacement used in~\cite[Lemma~3.2]{AEY2022}.
\end{proof}

\section{Constructions, endpoints, and threshold bounds}\label{sec:foundations}

\subsection{A common construction family}

We use the following consequence of Hall's theorem. A balanced bipartite
graph with $a$ vertices per side and minimum degree at least $a/2$ has a
perfect matching. Indeed, a nonempty set $S$ on one side with $|S|\le a/2$
has at least $a/2$ neighbors. If $|S|>a/2$, every vertex on the opposite
side has a neighbor in $S$, since otherwise its degree would be less than
$a/2$.

The recipe below is the family of~\cite[Proposition~4.1]{BET2023}.
We include its degree and matching verification and use the same Hall
argument to establish factor-criticality throughout the stated range.

\begin{lemma}[Intermediate construction]\label{input:construction}
Let $d\ge3$, $t\ge0$, and $2t+2\le d$. There is a triangle-free
factor-critical graph $B_{d,d+t}$ with
\[
  |V(B_{d,d+t})|=2(d+t)+1,\quad
  \nu(B_{d,d+t})=d+t,\quad
  e(B_{d,d+t})=d(d+t)+t+1,
\]
whose $2(d+t)$ original vertices have degree $d$ and whose remaining vertex
has degree $2t+2$.
\end{lemma}

\begin{proof}
Start with $K_{d+t,d+t}$ on sides $L,R$. Partition each side into a set $F$
of size $d-1$ and a set $H$ of size $t+1$. Delete $t$ edge-disjoint perfect
matchings between $F_L$ and $F_R$, delete all edges between $H_L$ and $H_R$,
and add a vertex $v$ adjacent to $H_L\cup H_R$. The required disjoint matchings
exist, for example by taking cyclic shifts between the two sets of size
$d-1$.

Vertices in $F$ have degree $d+t-t=d$, those in $H$ have degree
$d-1+1=d$, and $v$ has degree $2t+2$. Triangle-freeness follows because the
graph without $v$ is bipartite and the neighbors of $v$ are independent.
The edge count follows by summing degrees.

Our hypotheses imply $d\ge t+3$: this is immediate for $t=0$, and follows
from $2t+2\ge t+3$ for $t\ge1$. Deleting $v$ leaves a balanced bipartite
graph with $d+t$ vertices per side and minimum degree $d-1\ge(d+t)/2$.
If $x\in L$ is deleted instead, match $v$ to any vertex of $H_R$.
The remaining bipartite graph has $d+t-1$ vertices per side and minimum
degree at least $d-2\ge(d+t-1)/2$. Hall's theorem supplies a perfect
matching in both cases; deletion from $R$ is symmetric. Thus the graph
is factor-critical and has matching number $d+t$.
\end{proof}

\subsection{The left endpoint}

\begin{proposition}\label{input:left}
For $d\ge2$,
\[
  \fDelta(d,m)=dm\quad(0\le m<d),\qquad
  \fDelta(d,d)=d^2+1.
\]
\end{proposition}

\begin{proof}
If $m<d$, \cref{lem:component-lower} leaves only stars in a normalized
extremal graph. If $m=d$, there is at most one non-star component, and it
has matching number $d$. Its edge count is at most $d^2+1$ by
\cref{lem:fc-order-bound}. These bounds are attained by $m$ stars, and by
$B_{d,d}$ from \cref{input:construction}, respectively. For $d=2$, the latter
construction is replaced by $C_5$. These are the endpoint arguments
of~\cite[Theorems~3.3, 3.4]{AEY2022}.
\end{proof}

\subsection{The threshold and the upper limit on component size}

We retain one construction theorem from~\cite[Proposition~4.3]{AEY2022}:
for every $d\ge2$ there is a triangle-free factor-critical graph with the
threshold degree profile and matching number
\[
  U(d)=
  \begin{cases}
    \floor{5d/4},&d\text{ even},\\
    \floor{5(d+1)/4},&d\text{ odd}.
  \end{cases}
\]
These are the graphs denoted $B_d$ in that paper. We use their existence,
not a claim that they realize the least threshold in every degree.
This is the only problem-specific construction theorem used below whose
proof is not reproduced in this draft.
It follows that $Z(d)$ exists. Let $C_d$ be a graph attaining this least
threshold. Its degree profile gives
\[
  e(C_d)=dZ(d)+\floor{d/2}.
\]

\begin{proposition}\label{input:Z}
For all $d\ge2$,
\[
  d\le Z(d)\le U(d)<2d.
\]
Moreover, $Z(d)\ge d+1$ for $d\ge4$, $Z(d)=\floor{5d/4}$ for even $d$,
and $Z(3)=3$, $Z(5)=6$.
\end{proposition}

\begin{proof}
A threshold graph of matching number $j<d$ would have $dj+\floor{d/2}>dj$
edges, contrary to \cref{input:left}. If $d\ge4$, one of matching number
$d$ would have more than $d^2+1$ edges, also impossible. This proves the
lower bounds; existence of the cited construction gives $Z(d)\le U(d)<2d$.

For even $d$, apply the Andr\'asfai--Erd\H{o}s--S\'os theorem in the form
that every triangle-free graph on $n$ vertices with minimum degree greater
than $2n/5$ is bipartite; see~\cite{AES1974} and
\cite[Corollary~2.6]{AEY2022}.
The factor-critical threshold graph is nonbipartite and $d$-regular, so
$5d\le4Z(d)+2$. Since $d$ is even, this gives
$Z(d)\ge\floor{5d/4}$, as in~\cite[Lemma~4.7]{AEY2022}.
Finally, $B_{3,3}$ and $B_{5,6}$ from \cref{input:construction} have the
required almost-regular profiles, attaining the corresponding lower bounds.
\end{proof}

\begin{proposition}[Component reduction]\label{input:structure}
For every $d\ge2$ and $m\ge0$, an extremal graph can be chosen with each
component either a $d$-star or a factor-critical extremal graph of matching
number $j$ satisfying $d\le j\le Z(d)$.
\end{proposition}

\begin{proof}
Only the upper limit remains after \cref{lem:star-reduction,lem:component-lower}.
If a non-star component $H$ has matching number $j>Z(d)$, then
\[
  e(H)\le dj+\floor{d/2}.
\]
Replacing it by $C_d$ and $j-Z(d)$ stars attains that many edges and increases
the number of star components, contradicting normalization. This is the
argument of~\cite[Lemma~4.5]{AEY2022}.
\end{proof}

\begin{proposition}\label{input:threshold}
For every $d\ge2$,
\[
  \fDelta(d,Z(d))=dZ(d)+\floor{d/2}.
\]
\end{proposition}

\begin{proof}
The graph $C_d$ attains the stated value. Since $Z(d)<2d$, a normalized
extremal graph has at most one non-star component. If its matching number
is $j$, degree summation bounds its edges by $dj+\floor{d/2}$; the stars
contribute at most $d(Z(d)-j)$. If there is no such component, the bound is
immediate. This is the threshold case of~\cite[Theorem~4.8]{AEY2022}.
\end{proof}

\subsection{Deficiency notation}

For the new fixed-order argument, write
\[
  i=d+t,\qquad n=2i+1=2d+2t+1,\qquad
  L=di+t+1=d(d+t)+t+1.
\]
For any graph with maximum degree at most $d$, put
\[
  \epsilon(v)=d-\deg(v),\qquad D(G)=\sum_v\epsilon(v)=dn-2e(G).
\]
Then
\begin{equation}\label{eq:target-deficiency}
  e(G)\le L\quad\Longleftrightarrow\quad D(G)\ge d-2t-2.
\end{equation}
The deficiency is nonnegative and has the same parity as $d$, because $n$
is odd.

\section{A saturated five-cycle}\label{sec:saturated-cycle}

The following symmetrization supplies the saturated cycle needed by the structural argument.

\begin{lemma}[Saturated five-cycle]\label{lem:saturated-cycle}
Let $d\ge7$ and $1\le t<d$. Suppose that $G$ is a triangle-free graph on
$n=2d+2t+1$ vertices, with $\Delta(G)\le d$. If $e(G)>L$, then $G$ contains
an induced $5$-cycle all of whose vertices have degree $d$.
\end{lemma}

\begin{proof}
Since the edge count is integral, the counterexample assumption and
\cref{eq:target-deficiency} give
\begin{equation}\label{eq:counter-deficiency}
  D(G)\le d-2t-4.
\end{equation}
If $G$ is bipartite, its smaller part has at most $d+t$ vertices. Summing
degrees over this part gives $e(G)\le d(d+t)<L$, a contradiction. Thus $G$ is
not bipartite.

Let $C$ be a shortest odd cycle, of length $\ell$. It is induced. Every
vertex outside $C$ has at most two neighbors on $C$. Indeed, suppose that an
outside vertex has at least three such neighbors. The cyclic gaps between
successive neighbors have length at least two, because $G$ is triangle-free.
At least one gap is odd, and among at least three gaps such an odd gap has
length at most $\ell-4$. That gap together with the two incident edges from
the outside vertex gives a shorter odd cycle.

Every vertex of $G$ therefore sends at most two incidences into $C$, and hence
\[
  \sum_{v\in C}\deg(v)\le2n.
\]
It follows that
\[
  D(G)\ge\sum_{v\in C}(d-\deg(v))\ge \ell d-2n.
\]
If $\ell\ge7$, then
\[
  D(G)\ge3d-4t-2>d-2t-4,
\]
where the strict inequality is $2(d-t+1)>0$. This contradicts
\cref{eq:counter-deficiency}. Triangle-freeness excludes $\ell=3$, so $G$
contains a $5$-cycle.

Choose a $5$-cycle
\[
  C=c_0c_1c_2c_3c_4c_0
\]
maximizing $\sum_{j=0}^4\deg(c_j)$. Indices are read modulo five. Every
outside vertex has at most two neighbors on $C$. Let $R$ be the set of outside
vertices with fewer than two neighbors on $C$, and put
\[
  q=\sum_{x\in R}\bigl(2-|\N(x)\cap C|\bigr),
  \qquad r=|R|.
\]
Thus $r\le q$. Let $D_C$ and $D_{\mathrm{out}}$ be the sums of the
deficiencies on and outside $C$. Counting incidences with $C$ gives
\begin{equation}\label{eq:cycle-deficiency}
  D_C=5d-2n+q,
  \qquad
  q+D_{\mathrm{out}}=D(G)-5d+2n\le2t-2.
\end{equation}

For each $j$, let $A_j$ consist of $c_j$ and every outside vertex whose two
neighbors on $C$ are $c_{j-1}$ and $c_{j+1}$. These five sets partition
$V(G)\setminus R$. Write
\[
  a_j=|A_j|,
  \qquad
  b_j=|\N(c_j)\cap R|.
\]
The definition of the bags gives
\[
  \deg(c_j)=a_{j-1}+a_{j+1}+b_j.
\]
Adding this identity for $j+2$ and $j+3$ yields
\begin{align}
  a_j
    &=n-r-\deg(c_{j+2})-\deg(c_{j+3})+b_{j+2}+b_{j+3}\notag\\
    &\ge2t+1-r.\label{eq:bag-lower}
\end{align}
Consequently,
\[
  |A_j\setminus\{c_j\}|\ge2t-r>D_{\mathrm{out}},
\]
where the strict inequality follows from \cref{eq:cycle-deficiency} and
$q\ge r$. Every vertex of degree less than $d$ contributes at least one to
$D_{\mathrm{out}}$. Therefore $A_j\setminus\{c_j\}$ contains a vertex $x_j$
of degree $d$. Replacing $c_j$ by $x_j$ produces another $5$-cycle. The
maximal choice of $C$ now implies $\deg(c_j)=d$. This holds for every $j$.
\end{proof}

The bags $A_j$ are independent. Moreover, edges between distinct bags can
occur only between consecutive bags: vertices in any forbidden pair share a
neighbor on $C$, so an edge between them would create a triangle.

\section{Stability around a saturated five-cycle}\label{sec:hard-case}

The natural parameter is the \emph{cycle excess}
\[
  q=2n-5d.
\]
For a saturated five-cycle, this is exactly the number of missing incidences
into the cycle:
\[
  q=\sum_{x\in R}(2-|\N(x)\cap C|).
\]
In particular, $q\ge0$. The original bags already form a subgraph of a
blow-up of $C_5$; the question is whether the exceptional vertices can be
assigned to the same five positions.

We prove a structural statement for arbitrary order $n$, independent of the
matching-number problem. Put
\[
  T=8d-3n.
\]
This quantity is the deficiency forced by an eight-vertex subgraph with
independence number three: each graph vertex sends at most three incidences
into that subgraph.

\begin{theorem}[Saturated-cycle stability]\label{thm:stability}
Let $G$ be a triangle-free graph on $n$ vertices with $\Delta(G)\le d$,
containing a five-cycle $C=c_0c_1c_2c_3c_4c_0$ whose vertices all have
degree $d$. Set $D=dn-2e(G)$, $q=2n-5d$, and $T=8d-3n$.
If either
\[
  D<T,
  \qquad\text{or}\qquad
  D\le T,\quad T>0,\quad q\ge2,
\]
then there is a homomorphism $G\to C_5$ taking $c_j$ to $j$.
\end{theorem}

The saturated-cycle hypothesis is essential to the scope of this statement:
we do not claim here that every graph near the extremal edge count has such
a cycle. The theorem itself has examples, such as balanced complete blow-ups
of $C_5$, for which $D=q=0$ and $T>0$.

\subsection{The two local obstructions}

Two ways of failing to extend the bag assignment lead to the same graph $H$.
It has distinct vertices $a,b,u,v,p_1,p_2,x_1,x_2$, the four edges between
$\{a,b\}$ and $\{u,v\}$, and the two paths
\[
  a-p_1-p_2-b,\qquad u-x_1-x_2-v.
\]
Its independence number is three. An independent set cannot meet both
$\{a,b\}$ and $\{u,v\}$. If it meets their union once, the edges $p_1p_2$
and $x_1x_2$
allow at most two further vertices; if it avoids the union, they allow at
most two. If it contains both $a,b$, it excludes $p_1,p_2$ and contains at
most one of $x_1,x_2$; the case of $u,v$ is symmetric.

A triangle-free supergraph of $H$ on these eight vertices may add edges only
between $\{p_1,p_2\}$ and $\{x_1,x_2\}$, and the added edges form a matching.
Adding two gives the eight-vertex M\"obius ladder $M$, the graph on
$\Z/8\Z$ with differences $1$, $4$, or $7$ adjacent. Thus $M$ arises as the
unique two-edge augmentation, up to isomorphism, of the obstruction to
extending the coloring.

The graph $M$ is triangle-free and $3$-regular. Its independence number is
three, and its independent triples are precisely the neighborhoods
\[
  \{j-1,j+1,j+4\},\qquad j\in\Z/8\Z.
\]
Indeed, an independent four-set in the $8$-cycle must be alternating, and
both alternating sets contain an opposite chord. For an independent triple,
the cyclic gaps are at least two and sum to eight. The gap pattern $(2,2,4)$
contains an opposite chord; $(2,3,3)$ gives the displayed neighborhoods.
Distinct nonadjacent vertices of $M$ have a common neighbor.

\begin{lemma}[Obstruction cost and equality]\label{lem:no-mobius}
Every copy of $M$ in $G$ forces $D\ge T$. If $T>0$ and $q\ge2$, equality
$D=T$ is impossible.
\end{lemma}

\begin{proof}
Let $S$ be the eight vertices of a copy. Every $\N(x)\cap S$ is independent
in $M$, so has size at most three. Therefore
\[
  D_S:=\sum_{v\in S}(d-\deg_G(v))\ge8d-3n=T.
\]
If $D=T$, equality holds throughout: all deficiency lies on $S$, every
vertex outside $S$ has degree $d$, and every vertex has exactly three
neighbors in $S$. In particular, $G[S]=M$.

Partition $V(G)$ into bags $B_j$ according to the neighborhood
$\{j-1,j+1,j+4\}$ in $S$, with $j\in B_j$. Put $a_j=|B_j|$ and
$\delta_j=d-\deg_G(j)$. The common-neighbor property makes every bag
independent and permits edges between bags only along edges of $M$. Also
\[
  \deg_G(j)=a_{j-1}+a_{j+1}+a_{j+4}.
\]
Adding this identity at $j-1,j,j+1$ counts every bag once and $B_j$ twice,
so
\[
  a_j=3d-n-(\delta_{j-1}+\delta_j+\delta_{j+1})
  \ge3d-n-D_S=2n-5d=q.
\]
If $q\ge2$, each bag contains a vertex outside $S$. That vertex is saturated
and has no more possible neighbors than $j$, forcing $\deg_G(j)=d$.
Thus all of $S$ is saturated, contradicting $D_S=T>0$.
\end{proof}

\begin{lemma}\label{lem:no-H}
If $G$ contains no $M$, every copy of $H$ forces $D\ge T+2$.
\end{lemma}

\begin{proof}
On the vertex set $S$ of a copy of $H$, at most one edge can be added:
two added edges would give $M$. Hence $e(G[S])\le11$. Every outside
vertex has at most three neighbors in $S$, giving
\[
  \sum_{v\in S}\deg_G(v)\le22+3(n-8)=3n-2.
\]
Consequently $D\ge D_S\ge8d-3n+2=T+2$.
\end{proof}

\subsection{Extending the cycle coloring}

\begin{lemma}[Extension criterion]\label{lem:extension}
If $G$ contains a saturated five-cycle, contains no $H$, and $D\le T$,
then its natural bag assignment extends to a homomorphism $G\to C_5$.
\end{lemma}

\begin{proof}
Use the bags $A_j$ and exceptional set $R$ of
\cref{sec:saturated-cycle}, and write $r=|R|$ and
$b_j=|\N(c_j)\cap R|$. Saturation gives
\[
  r\le q,\qquad \sum_jb_j=2r-q.
\]
The exact bag identity therefore yields the uniform bound
\begin{equation}\label{eq:hard-bag-size}
  a_j=n-r-2d+b_{j+2}+b_{j+3}\le n-2d.
\end{equation}
Indeed, $b_{j+2}+b_{j+3}\le2r-q\le r$.

No $x\in R$ can have neighbors in consecutive bags. Otherwise, after
relabeling, choose $u\in A_0$ and $v\in A_1$ adjacent to $x$. There is a
copy of $H$ with
\[
  (a,b,u_H,v_H,p_1,p_2,x_1,x_2)
  =(c_0,c_2,v,c_1,c_4,c_3,x,u).
\]
The two paths are $c_0-c_4-c_3-c_2$ and $v-x-u-c_1$, and the four cross-edges
come from the bag definitions. The vertices are distinct: $u=c_0$ or
$v=c_1$ would create a triangle with $x$.

For every $x\in R$, the deficiency bound gives
\begin{equation}\label{eq:R-neighbors}
  |\N(x)\setminus R|
  \ge d-D-r+1
  \ge d-T-q+1=n-2d+1.
\end{equation}
This exceeds the size of any one bag. Its occupied positions form an
independent set in $C_5$, so $x$ meets exactly two nonconsecutive bags.
Assign $x$ to their unique common neighboring position.

Only edges inside $R$ remain to check. If an edge $xy$ has equal or
nonadjacent assigned positions, the two occupied bag-pairs share a bag
$A_j$. Choose $u\in\N(x)\cap A_j$ and $v\in\N(y)\cap A_j$.
Triangle-freeness gives $u\ne v$. The vertices $x,y,u,v$ and the four cycle
vertices other than $c_j$ contain $H$, with
\[
  (a,b,u_H,v_H,p_1,p_2,x_1,x_2)
  =(c_{j-1},c_{j+1},u,v,c_{j-2},c_{j+2},x,y).
\]
They are distinct even if $u$ or $v$ equals $c_j$, since $c_j$ was omitted.
This contradiction proves that every edge maps to an edge of $C_5$.
\end{proof}

\begin{proof}[Proof of \cref{thm:stability}]
Under either hypothesis, \cref{lem:no-mobius} excludes $M$.
Since $D\le T$, \cref{lem:no-H} then excludes $H$, and
\cref{lem:extension} extends the coloring.
\end{proof}

\section{Odd excess costs deficiency}\label{sec:capacity}

Once the coloring is available, the excess has a second interpretation.
Let $B_0,\ldots,B_4$ be its color classes, each containing the corresponding
saturated cycle vertex, and set
\[
  s_j=|B_j|,\qquad
  p_j=s_{j-1}+s_{j+1}-d,\qquad h=n-2d.
\]
Saturation ensures $p_j\ge0$. Direct summation gives
\begin{equation}\label{eq:potential-identities}
  \sum_jp_j=q,\qquad s_j=h-p_{j+2}-p_{j+3}.
\end{equation}
Thus the same excess $q$ counts both the original missing cycle incidences
and the surplus neighboring capacity after coloring.

Let $F$ be the complete blow-up on these classes, put
$\mu=e(F)-e(G)$, and write $w_j=s_jp_j$ and $W=\sum_jw_j$.
Every vertex in $B_j$ has $d+p_j$ possible neighbors, so
\begin{equation}\label{eq:missing-deficiency}
  D=2\mu-W.
\end{equation}
For any independent set $I$ of cycle positions, a missing edge has at most
one endpoint in its classes. Since at least $w_j$ incidences are missing
from $B_j$,
\begin{equation}\label{eq:independent-potential}
  D\ge2\sum_{j\in I}w_j-W.
\end{equation}

The following argument replaces the capacity table. It applies more
generally whenever the positions with positive excess have bipartite support.

\begin{lemma}[Odd excess imbalance]\label{lem:odd-excess}
Suppose $p_j$ are nonnegative integers with odd sum $q=2k-1>0$, and
the subgraph of $C_5$ on $\{j:p_j>0\}$ is bipartite. If $h\ge q$ and
$w_j=p_j(h-p_{j+2}-p_{j+3})$, then some independent set $I$ satisfies
\[
  2\sum_{j\in I}w_j-\sum_jw_j\ge h-k(k-1).
\]
In particular, for $q=1$ or $q=3$, \cref{eq:independent-potential} gives
$D\ge h$ or $D\ge h-2$, respectively.
\end{lemma}

\begin{proof}
Bipartition the positive support into independent sets $I,J$, choosing
$I$ to have larger total $p$-weight. Put $P=\sum_{j\in I}p_j$; integrality
and oddness give $k\le P\le q$.

Expand the difference of the two $w$-weights. Products $p_ip_j$ from
nonadjacent positions on opposite sides cancel. Those within $J$ have
positive sign and may be discarded. Those within $I$ cost at most
\[
  2\sum_{\{i,j\}\subseteq I}p_ip_j
  =P^2-\sum_{i\in I}p_i^2\le P(P-1),
\]
because the $p_i$ are integers. Therefore
\[
  2\sum_{j\in I}w_j-W
  \ge h(2P-q)-P(P-1)\ge h-k(k-1).
\]
The last inequality follows from
\[
  h(2P-q)-P(P-1)-\bigl(h-k(k-1)\bigr)
  =(P-k)(2h-P-k+1)\ge0,
\]
using $P\ge k$ and $P,k\le q\le h$.
For $q\le3$, the positive support has at most three positions, so it is
automatically bipartite.
\end{proof}

\section{The fixed-order bound}\label{sec:direct-cases}

We now apply the structural and capacity statements with
$n=2d+2t+1$. The relevant parameters are
\begin{equation}\label{eq:excess-parameters}
  q=4t+2-d,\qquad h=2t+1,\qquad
  T=2d-6t-3.
\end{equation}

\begin{proposition}\label{prop:excess-bound}
Let $d\ge7$ and $1\le t<d$. If $q=4t+2-d$ is negative, one, or three,
then every triangle-free graph on $2d+2t+1$ vertices with maximum degree
at most $d$ has at most $L=d(d+t)+t+1$ edges.
\end{proposition}

\begin{proof}
A counterexample has a saturated five-cycle by \cref{lem:saturated-cycle}.
Hence $q\ge0$, disposing of negative excess. Moreover,
\begin{equation}\label{eq:excess-budget}
  D\le d-2t-4=T+q-3.
\end{equation}
If $q=1$, then $D\le T-2<T$. If $q=3$, then $d=4t-1$ and
$D\le T=2t-5$. Nonnegativity of $D$ forces $t\ge3$, so $T>0$.
In either case \cref{thm:stability} gives a homomorphism to $C_5$
preserving the saturated cycle. Since $h=2t+1\ge q$, the odd excess
imbalance gives, in both cases,
\[
  D\ge h-(q-1)=2t+2-q.
\]
But \cref{eq:excess-budget} also reads $D\le2t-2-q$, a contradiction.
\end{proof}

\section{Determining the threshold}\label{sec:exact-threshold}

We identify the least threshold before using it in the general formula.
The following argument applies the fixed-order proposition directly, without
assuming the desired value of $Z(d)$.

\begin{proposition}\label{cor:odd-Z}
For odd $d\ge9$,
\[
  Z(d)=\floor{\frac{5(d+1)}4}.
\]
\end{proposition}

\begin{proof}
The upper bound is \cref{input:Z}. Suppose a threshold graph occurs at
matching number $j<U(d)$. The same proposition gives $j\ge d+1$.
Write $j=d+t$. Since $j<U(d)$, we have $4t\le d+1$ and $t<d$.
The odd integer $q=4t+2-d$ is at most three, so it is negative, one, or
three. By \cref{prop:excess-bound}, the graph has deficiency at least
\[
  d-2t-2\ge(d-5)/2\ge2.
\]
Its threshold degree profile has deficiency one, a contradiction.
\end{proof}

\begin{lemma}\label{lem:seven-threshold}
We have $Z(7)=9$.
\end{lemma}

\begin{proof}
The construction $B_{7,9}$ from \cref{input:construction}, with $t=2$,
is triangle-free and factor-critical, with eighteen vertices of degree
seven and one of degree six. It therefore gives $Z(7)\le9$ directly.

For the reverse bound, \cref{input:Z} gives $Z(7)\ge8$. A threshold graph
of matching number eight would have seventeen vertices and
$(7\cdot17-1)/2=59$ edges. But \cref{prop:excess-bound}, with $d=7,t=1$,
allows at most $7\cdot8+2=58$. Thus $Z(7)=9$.
\end{proof}

\begin{proposition}[Explicit threshold]\label{prop:explicit-Z}
For every $d\ge2$, $Z(d)=z_d$, where $z_d$ is defined in
\cref{eq:explicit-threshold}.
\end{proposition}

\begin{proof}
Combine the even-degree and small odd values of \cref{input:Z} with
\cref{cor:odd-Z,lem:seven-threshold}.
\end{proof}

\section{The intermediate component bound}\label{sec:conj21}

\begin{theorem}\label{thm:fixed-order}
Let $d\ge7$ and $d<i<Z(d)$. Every triangle-free graph on $2i+1$ vertices
with maximum degree at most $d$ has at most $di+i-d+1$ edges.
\end{theorem}

\begin{proof}
Write $i=d+t$; \cref{input:Z} gives $1\le t<d$.
If $d$ is even, $i<\floor{5d/4}$ implies $4i\le5d-4$, so
$q=4i+2-5d\le-2$.
If $d$ is odd, $i<\floor{5(d+1)/4}$ implies $4i\le5d+1$,
so $q\le3$ and $q$ is odd.
In either case \cref{prop:excess-bound} applies.
\end{proof}

\begin{proposition}[Intermediate value]\label{thm:conj21}
For $7\le d<i<Z(d)$,
\[
  \fDelta(d,i)=di+i-d+1.
\]
\end{proposition}

\begin{proof}
Write $i=d+t$ and use the normalized graph of \cref{input:structure}.
Since $i<2d$, there is at most one non-star component. If there is none,
there are at most $di$ edges. Otherwise let its matching number be
$j=d+s\le i$. For $s=0$, \cref{input:left} gives at most $dj+1$ edges;
for $s>0$, \cref{thm:fixed-order} gives at most $dj+s+1$.
Together with the stars this is at most
\[
  di+s+1\le di+t+1.
\]

For attainment, the bounds $i<U(d)$ imply $2t+2\le d$ when $d\ge7$.
Explicitly, $4t\le d-4$ for even $d$, and $4t\le d+1$ for odd $d$.
Thus \cref{input:construction} supplies $B_{d,d+t}$ with the stated value.
\end{proof}

\begin{remark}
The fixed-order theorem requires no factor-criticality assumption. The
component reduction uses this stronger conclusion only on odd-order
factor-critical components.
\end{remark}

\section{The general formula as a packing theorem}\label{sec:packing}

For $d\ge4$, the component bounds above reduce all matching bounds to one
packing problem. Set
\[
  z=Z(d),
  \qquad
  b=\floor{d/2}.
\]
The explicit threshold values in \cref{prop:explicit-Z} imply
\begin{equation}\label{eq:threshold-bounds}
  d<z<2d,
  \qquad
  z\le d+b-1.
\end{equation}

Relative to the baseline $dm$, a component of matching number $j$ has utility
\[
  u(j)=
  \begin{cases}
    j-d+1,&d\le j<z,\\
    b,&j=z.
  \end{cases}
\]
At $j=d$, the first line follows from \cref{input:left}; for $d<j<z$,
it follows from \cref{thm:conj21}. This interval is empty for $d=4,5,6$,
since then $z=d+1$. The second line follows from \cref{input:threshold}.
Stars have zero utility and fill unused matching capacity. This reduction
is the component-packing viewpoint of~\cite[Lemma~4.9]{AEY2022}; the
normalization below proves the needed optimization directly.

\begin{lemma}[Packing normalization]\label{lem:packing}
Any feasible list of component weights in $[d,z]$ can be replaced, without
increasing its total weight or decreasing its utility, by a list containing as
many weight-$z$ components as the capacity permits and at most one component
of weight below $z$.
\end{lemma}

\begin{proof}
Take two weights $a,c<z$, and put $s=a+c$. Since $s\ge2d>z$, replace them by
one weight-$z$ component. If $s-z\ge d$, add a component of weight $s-z$;
otherwise fill the remaining capacity with stars. The total weight does not
increase.

If $s-z\ge d$, the utility gain is
\[
  b+(s-z-d+1)-(s-2d+2)=b+d-z-1\ge0
\]
by \cref{eq:threshold-bounds}. If $s-z<d$, then $s\le z+d-1$ and the old
utility is
\[
  s-2d+2\le z-d+1\le b.
\]
Thus the replacement never decreases utility and strictly decreases the
number of sub-threshold components. Repetition leaves at most one.

Now write the total capacity as $m=kz+r$, with $0\le r<z$, and suppose the
normalized list contains $k'<k$ weight-$z$ components. Discard its one smaller
component if present and add one weight-$z$ component. The new weight is
$(k'+1)z\le kz\le m$, and the utility does not decrease because $u(j)\le b$
for every $j<z$. Repeating gives exactly $k$ threshold components. The
remaining capacity is $r$, in which at most one residual component can be
placed.
\end{proof}

\begin{proof}[Proof of \cref{thm:main}]
First \cref{prop:explicit-Z} identifies $Z(d)=z_d$.

For $d=2,3$, this threshold equals $d$. By \cref{input:structure}, every
non-star component of a normalized extremal graph has matching number $d$,
and contributes at most $d^2+\floor{d/2}$ edges. Hence at most
$\floor{m/d}$ such components can occur, giving
\[
  \fDelta(d,m)\le dm+\floor{m/d}\floor{d/2}.
\]
This is attained by that many copies of $C_d$ and the remaining stars.
It agrees with \cref{eq:main-formula}, since $r<z_d=d$.

Now let $d\ge4$ and $m=kz_d+r$. The component reduction and
\cref{lem:packing} give $k$ threshold components and at most one residual
non-star component in an optimal packing. If $r<d$, the residual utility is
zero. If $r\ge d$, it is at most $r-d+1$. Therefore
\[
  \fDelta(d,m)\le dm+k\floor{d/2}+\max\{0,r-d+1\}.
\]

For attainment, take $k$ copies of $C_d$. If $r<d$, add $r$ stars.
If $r=d$, add the left-endpoint graph from \cref{input:left}.
If $r>d$, use $B_{d,r}$ from \cref{input:construction}. This last case
can occur only for $d\ge7$, and the parameter check in
\cref{thm:conj21} applies. These graphs have matching number $m$ and attain
the upper bound, completing the proof for every $m\ge0$.
\end{proof}

\section{Consequences}\label{sec:consequences}

\begin{corollary}[The two conjectures]\label{thm:conj22}
Conjectures~2.1 and~2.2 of~\cite{BET2023} hold: the intermediate value is
$\fDelta(d,i)=di+i-d+1$ for $7\le d<i<Z(d)$, and the general formula in
terms of $Z(d)$ holds for every $d\ge2$ and $m\ge0$.
\end{corollary}

\begin{proof}
The first assertion is \cref{thm:conj21}. The second is \cref{thm:main}
with $z_d=Z(d)$ from \cref{prop:explicit-Z}.
\end{proof}

\begin{corollary}\label{cor:factor-critical}
Let $d\ge2$ and $d<m<Z(d)$. Every extremal triangle-free graph $G$ with
$\Delta(G)\le d$, $\nu(G)=m$, and $|V(G)|=2m+1$ is factor-critical.
\end{corollary}

\begin{proof}
If $G$ is not factor-critical, there is a vertex $v$ such that $G-v$ has no
perfect matching. Since $G-v$ has $2m$ vertices, its matching number is at most
$m-1$. Therefore
\[
  e(G)\le\fDelta(d,m-1)+d=\fDelta(d,m)-1,
\]
where the equality follows from \cref{thm:conj21}, including the known
left endpoint when $m=d+1$. This contradicts extremality.
\end{proof}

\section{Verification status and points for review}\label{sec:verification}

The argument above is a proof draft and has not yet been independently
refereed. Lean 4.33.0 with Mathlib checks the following parts:

\begin{itemize}
  \item the graph argument from shortest odd cycles through the saturated
  $5$-cycle, the $M$ and $H$ exclusions, the $C_5$-coloring, and every
  fixed-order degree case;
  \item the required capacity inequalities, the general stability theorem,
  and the terminating packing normalization;
  \item boundedness and attainment for $\fDelta(d,m)$ and $Z(d)$, the full
  Proposition~4.1 family $B_{d,d+t}$ (order, triangle-freeness, degrees,
  matching number, and edge count), and the factor-critical graph $B_{7,9}$
  used for the degree-seven threshold.
\end{itemize}

The integrated mathematical proof above includes the short published
component and endpoint arguments. The Lean development still represents
those results, the general threshold constructions, the disjoint-union lower
construction for the full formula, and the previously
solved small cases through the explicit interface
\texttt{AuditedLiteratureInputs}. Their inclusion as prose proofs here does
not discharge those formal hypotheses. From that interface, Lean proves the
fixed-order upper bound and the packing deduction, and derives both
conjectures in terms of $Z(d)$. Proposition~4.1 is no longer assumed by that
interface: Lean constructs its graph and derives the residual lower bounds
directly; the remaining disjoint-union field takes those bounds as premises.
The complete identification $Z(d)=z_d$ in
\cref{prop:explicit-Z} is proved in the manuscript. Lean additionally proves
the degree-seven case and \cref{cor:odd-Z} for every odd $d\ge9$, the latter
from the published bounds $d+1\le Z(d)\le\floor{5(d+1)/4}$. The general Hall proof of factor-criticality in
\cref{input:construction} remains a written proof.

The complete Lean build audits the named theorems listed in
\texttt{Audit.lean}. It contains no admitted proof, project-specific axiom,
or native-evaluation shortcut. The reported dependencies are only Lean's
standard foundational principles
(``propext'', classical choice, and quotient soundness). Lean formalizes the
weighted bipartition proof of \cref{lem:odd-excess} and uses it for the
excess-one and excess-three capacity conclusions. The earlier finite pattern
proofs remain as independent regression checks.
For $B_{7,9}$, Lean checks explicit matching partners after each of the
nineteen vertex deletions, in place of the Hall argument given above.
The upper bound $Z(7)\le9$ has no external hypotheses; its lower bound uses
only the general published lower bound $8\le Z(7)$.
The formal files and the exact map from assumptions to published theorem
numbers will accompany a later version of this draft.

The main points for mathematical review are the transition from normalized
components to the fixed-order problem, \cref{lem:saturated-cycle}, the two
obstruction embeddings in \cref{sec:hard-case}, the odd excess imbalance,
the threshold identification in \cref{sec:exact-threshold}, and the packing
replacement in \cref{lem:packing}. The attribution of the integrated
published arguments is recorded where each is used.

\begin{remark}[A discrepancy in the 2023 preprint]
The introductory discussion of~\cite{BET2023} gives $Z(9)=13$ once, whereas
its Section~4 and computational table give $Z(9)=12$. The latter is consistent
with the proven upper bound
$Z(9)\le\floor{5(9+1)/4}=12$ from~\cite{AEY2022}. No argument here uses the
value $13$. There is a second table inconsistency: Table~5 reports 148 edges
for $(d,m)=(8,18)$, while Table~6 reports 149. The latter agrees with the
conjectured formula and with the displayed decomposition into one
$Z(8)$-component and one component of matching number eight.
\end{remark}

\section*{Verification files}

The Lean development, the finite verification script, and the map from Lean
inputs to published results are available at
\url{https://github.com/erdembanak/extremal-IP/tree/master/proof}.

\begin{thebibliography}{9}

\bibitem{AES1974}
B.~Andr\'asfai, P.~Erd\H{o}s, and V.~T.~S\'os,
\newblock On the connection between chromatic number, maximal clique and
minimal degree of a graph,
\newblock \emph{Discrete Mathematics} 8 (1974), 205--218.

\bibitem{AEY2022}
M.~Ahanjideh, T.~Ekim, and M.~A.~Y\i ld\i z,
\newblock Maximum size of a triangle-free graph with bounded maximum degree
and matching number,
\newblock \emph{arXiv preprint arXiv:2207.02271}, 2022.
\newblock \url{https://arxiv.org/abs/2207.02271}.

\bibitem{BET2023}
A.~E.~Banak, T.~Ekim, and Z.~C.~Ta\c{s}k\i n,
\newblock Constructing extremal triangle-free graphs using integer
programming,
\newblock \emph{arXiv preprint arXiv:2304.01729}, 2023.
\newblock \url{https://arxiv.org/abs/2304.01729}.

\end{thebibliography}

\end{document}
